\begin{thebibliography}{99}
\bibitem{aristoff-lelievre-14}D. Aristoff \& T. Lelièvre - “Mathematical analysis of temperature accelerated dynamics”, Multiscale Model. Simul. 12 (2014) no. 1, p. 290-317 | MR | DOI | Zbl
\bibitem{bai2010efficient}X. Bai, A. F. Voter, R. G. Hoagland, M. Nastasi \& B. P. Uberuaga - “Efficient annealing of radiation damage near grain boundaries via interstitial emission”, Science 327 (2010) no. 5973, p. 1631-1634 | DOI
\bibitem{Ber}N. Berglund - “Kramers’ law: validity, derivations and generalisations”, Markov Process. Related Fields 19 (2013) no. 3, p. 459-490 | MR | Zbl
\bibitem{BD15}N. Berglund \& S. Dutercq - “The Eyring-Kramers law for Markovian jump processes with symmetries”, J. Theoret. Probab. 29 (2016) no. 4, p. 1240-1279 | DOI | Zbl
\bibitem{Berglund_Gentz_MPRF}N. Berglund \& B. Gentz - “The Eyring-Kramers law for potentials with nonquadratic saddles”, Markov Process. Related Fields 16 (2010) no. 3, p. 549-598 | MR | Zbl
\bibitem{bianchi2020soft}A. Bianchi, A. Gaudillière \& P. Milanesi - “On soft capacities, quasi-stationary distributions and the pathwise approach to metastability”, J. Statist. Phys. 181 (2020) no. 3, p. 1052-1086 | MR | DOI | Zbl
\bibitem{borisov2}D. Borisov \& O. Sultanov - “Complete asymptotics for solution of singularly perturbed dynamical systems with single well potential”, Mathematics 8 (2020) no. 6, p. 949 | DOI
\bibitem{BEGK}A. Bovier, M. Eckhoff, V. Gayrard \& M. Klein - “Metastability in reversible diffusion processes. I. Sharp asymptotics for capacities and exit times”, J. Eur. Math. Soc. (JEMS) 6 (2004) no. 4, p. 399-424 | DOI | Zbl | MR
\bibitem{BGK}A. Bovier, V. Gayrard \& M. Klein - “Metastability in reversible diffusion processes. II. Precise asymptotics for small eigenvalues”, J. Eur. Math. Soc. (JEMS) 7 (2005) no. 1, p. 69-99 | DOI | MR | Zbl
\bibitem{bovier2016metastability}A. Bovier \& F. den Hollander - Metastability. A potential-theoretic approach, Grundlehren Math. Wissen., vol. 351, Springer, Cham, 2015 | DOI | Zbl | MR
\bibitem{brocker1982introduction}T. Bröcker \& K. Jänich - Introduction to differential topology, Cambridge University Press, Cambridge, 1982 | Zbl | MR
\bibitem{brown-94}R. Brown - “The mixed problem for Laplace’s equation in a class of Lipschitz domains”, Comm. Partial Differential Equations 19 (1994) no. 7-8, p. 1217-1233 | DOI | MR | Zbl
\bibitem{cameron-14b}M. Cameron - “Computing the asymptotic spectrum for networks representing energy landscapes using the minimum spanning tree”, Netw. Heterog. Media 9 (2014) no. 3, p. 383-416 | DOI | MR | Zbl
\bibitem{cattiaux-collet-lambert-martinez-meleard-san-martin-09}P. Cattiaux, P. Collet, A. Lambert, S. Martínez, S. Méléard \& J. San Martín - “Quasi-stationary distributions and diffusion models in population dynamics”, Ann. Probab. 37 (2009) no. 5, p. 1926-1969 | DOI | MR | Zbl
\bibitem{ChampagnatCoulibaly}N. Champagnat, K. A. Coulibaly-Pasquier \& D. Villemonais - “Criteria for exponential convergence to quasi-stationary distributions and applications to multi-dimensional diffusions”, in Séminaire de Probabilités XLIX, Lect. Notes in Math., vol. 2215, Springer, Cham, 2018, p. 165-182 | DOI | Zbl
\bibitem{CL}K. C. Chang \& J. Liu - “A cohomology complex for manifolds with boundary”, Topol. Methods Nonlinear Anal. 5 (1995) no. 2, p. 325-340 | DOI | MR | Zbl
\bibitem{ChPi}J. Chazarain \& A. Piriou - Introduction to the theory of linear partial differential equations, Studies in Math. and its Appl., vol. 14, North-Holland Publishing Co., Amsterdam-New York, 1982 | Zbl
\bibitem{collet-martinez-san-martin-13}P. Collet, S. Martínez \& J. San Martín - Quasi-stationary distributions. Markov chains, diffusions and dynamical systems, Probability and its Appl., Springer, Heidelberg, 2013 | DOI | Zbl | MR
\bibitem{CFKS}H. L. Cycon, R. G. Froese, W. Kirsch \& B. Simon - Schrödinger operators with application to quantum mechanics and global geometry, Texts and Monographs in Physics, Springer-Verlag, Berlin, 1987 | Zbl | MR
\bibitem{Day4}M. V. Day - “Exponential leveling for stochastically perturbed dynamical systems”, SIAM J. Math. Anal. 13 (1982) no. 4, p. 532-540 | MR | Zbl | DOI
\bibitem{Day2}M. V. Day - “On the exponential exit law in the small parameter exit problem”, Stochastics 8 (1983) no. 4, p. 297-323 | MR | Zbl | DOI
\bibitem{day1984a}M. V. Day - “On the asymptotic relation between equilibrium density and exit measure in the exit problem”, Stochastics 12 (1984) no. 3-4, p. 303-330 | DOI | MR | Zbl
\bibitem{day1987r}M. V. Day - “Recent progress on the small parameter exit problem”, Stochastics 20 (1987) no. 2, p. 121-150 | DOI | MR | Zbl
\bibitem{Day1}M. V. Day - “Large deviations results for the exit problem with characteristic boundary”, J. Math. Anal. Appl. 147 (1990) no. 1, p. 134-153 | MR | DOI | Zbl
\bibitem{di2017small}G. Di Gesù \& D. Le Peutrec - “Small noise spectral gap asymptotics for a large system of nonlinear diffusions”, J. Spectral Theory 7 (2017) no. 4, p. 939-984 | Zbl | DOI | MR
\bibitem{di-gesu-lelievre-le-peutrec-nectoux-17}G. Di Gesù, T. Lelièvre, D. Le Peutrec \& B. Nectoux - “Jump Markov models and transition state theory: the quasi-stationary distribution approach”, Faraday Discussions 195 (2017), p. 469-495 | DOI
\bibitem{DLLN}G. Di Gesù, T. Lelièvre, D. Le Peutrec \& B. Nectoux - “Sharp asymptotics of the first exit point density”, Ann. PDE 5 (2019) no. 1, article ID 5, 174 pages | DOI | MR | Zbl
\bibitem{DLLN-saddle1}G. Di Gesù, T. Lelièvre, D. Le Peutrec \& B. Nectoux - “The exit from a metastable state: concentration of the exit point distribution on the low energy saddle points, part 1”, J. Math. Pures Appl. (9) 138 (2020), p. 242-306 | Zbl | DOI | MR
\bibitem{DiSj}M. Dimassi \& J. Sjöstrand - Spectral asymptotics in the semi-classical limit, London Math. Soc. Lect. Note Series, vol. 268, Cambridge University Press, Cambridge, 1999 | DOI | Zbl | MR
\bibitem{eizenberg-90}A. Eizenberg - “The exponential leveling and the Ventcel-Freidlin ‘minimal action’ function”, J. Analyse Math. 54 (1990) no. 1, p. 99-111 | DOI | MR | Zbl
\bibitem{Eva}L. C. Evans - Partial differential equations, Graduate Studies in Math., vol. 19, American Mathematical Society, Providence, RI, 2010 | DOI | Zbl
\bibitem{evans-gariepy-18}L. C. Evans \& R. F. Gariepy - Measure theory and fine properties of functions, Textbooks in Math., CRC Press, Boca Raton, FL, 2015 | Zbl | DOI | MR
\bibitem{fan-yip-yildiz-14}Y. Fan, S. Yip \& B. Yildiz - “Autonomous basin climbing method with sampling of multiple transition pathways: application to anisotropic diffusion of point defects in hcp Zr”, J. Phys.: Condensed Matter 26 (2014), p. 365402
\bibitem{FrWe}M. I. Freidlin \& A. D. Wentzell - Random perturbations of dynamical systems, Grundlehren Math. Wissen., vol. 260, Springer, Heidelberg, 2012 | DOI | Zbl | MR
\bibitem{MR1814364}D. Gilbarg \& N. S. Trudinger - Elliptic partial differential equations of second order, Classics in Math., Springer-Verlag, Berlin, 2001 | Zbl | DOI
\bibitem{goldshtein-mitrea-mitrea-11}V. Gol’dshtein, I. Mitrea \& M. Mitrea - “Hodge decompositions with mixed boundary conditions and applications to partial differential equations on Lipschitz manifolds”, J. Math. Sci. 172 (2011) no. 3, p. 347-400 | DOI | Zbl
\bibitem{guillemin2010differential}V. Guillemin \& A. Pollack - Differential topology, vol. 370, American Mathematical Society, Providence, RI, 2010 | DOI | Zbl | MR
\bibitem{hanggi-talkner-barkovec-90}P. Hänggi, P. Talkner \& M. Borkovec - “Reaction-rate theory: fifty years after Kramers”, Rev. Modern Phys. 62 (1990) no. 2, p. 251-342 | DOI | MR
\bibitem{helffer-88}B. Helffer - Semi-classical analysis for the Schrödinger operator and applications, Lect. Notes in Math., vol. 1336, Springer-Verlag, Berlin, 1988 | DOI | Zbl | MR
\bibitem{helffer2013spectral}B. Helffer - Spectral theory and its applications, Cambridge Studies in Advanced Math., vol. 139, Cambridge University Press, Cambridge, 2013 | Zbl | DOI | MR
\bibitem{HKN}B. Helffer, M. Klein \& F. Nier - “Quantitative analysis of metastability in reversible diffusion processes via a Witten complex approach”, Mat. Contemp. 26 (2004), p. 41-85 | Zbl | MR
\bibitem{helffer-nier-06}B. Helffer \& F. Nier - Quantitative analysis of metastability in reversible diffusion processes via a Witten complex approach: the case with boundary, Mém. Soc. Math. France (N.S.), vol. 105, Société Mathématique de France, Paris, 2006 | Zbl | Numdam
\bibitem{HelSjII}B. Helffer \& J. Sjöstrand - “Puits multiples en limite semi-classique. II. Interaction moléculaire. Symétries. Perturbation”, Ann. Inst. H. Poincaré Phys. Théor. 42 (1985) no. 2, p. 127-212 | MR | Numdam | Zbl
\bibitem{helffer-sjostrand-85}B. Helffer \& J. Sjöstrand - “Puits multiples en mécanique semi-classique, IV: Etude du complexe de Witten”, Comm. Partial Differential Equations 10 (1985) no. 3, p. 245-340 | DOI | MR | Zbl
\bibitem{herau-hitrick-sjostrand-11}F. Hérau, M. Hitrik \& J. Sjöstrand - “Tunnel effect and symmetries for Kramers-Fokker-Planck type operators”, J. Inst. Math. Jussieu 10 (2011) no. 3, p. 567-634 | DOI | MR | Zbl
\bibitem{holley-kusuoka-stroock-89}R. A. Holley, S. Kusuoka \& D. W. Stroock - “Asymptotics of the spectral gap with applications to the theory of simulated annealing”, J. Funct. Anal. 83 (1989) no. 2, p. 333-347 | DOI | MR | Zbl
\bibitem{souga1}H. Ishii \& P. E. Souganidis - “Metastability for parabolic equations with drift: Part I”, Indiana Univ. Math. J. 64 (2015), p. 875-913 | Zbl | DOI | MR
\bibitem{souga2}H. Ishii \& P. E. Souganidis - “Metastability for parabolic equations with drift: part II. The quasilinear case”, Indiana Univ. Math. J. 66 (2017), p. 315-360 | Zbl | DOI | MR
\bibitem{jakab-mitrea-mitrea-09}T. Jakab, I. Mitrea \& M. Mitrea - “On the regularity of differential forms satisfying mixed boundary conditions in a class of Lipschitz domains”, Indiana Univ. Math. J. 58 (2009) no. 5, p. 2043-2071 | DOI | MR | Zbl
\bibitem{jost2008riemannian}J. Jost - Riemannian geometry and geometric analysis, Universitext, Springer-Verlag, Berlin, 2008 | DOI | Zbl | MR
\bibitem{kamin-78}S. Kamin - “Elliptic perturbation of a first order operator with a singular point of attracting type”, Indiana Univ. Math. J. 27 (1978) no. 6, p. 935-952 | DOI | MR | Zbl
\bibitem{kamin1979elliptic}S. Kamin - “On elliptic singular perturbation problems with turning points”, SIAM J. Math. Anal. 10 (1979) no. 3, p. 447-455 | DOI | MR | Zbl
\bibitem{landim2017dirichlet}C. Landim, M. Mariani \& I. Seo - “Dirichlet’s and Thomson’s Principles for Non-selfadjoint Elliptic Operators with Application to Non-reversible Metastable Diffusion Processes”, Arch. Rational Mech. Anal. 231 (2019) no. 2, p. 887-938 | DOI | MR | Zbl
\bibitem{laudenbach2011morse}F. Laudenbach - “A Morse complex on manifolds with boundary”, Geom. Dedicata 153 (2011) no. 1, p. 47-57 | DOI | MR | Zbl
\bibitem{le2012mathematical}C. Le Bris, T. Lelièvre, M. Luskin \& D. Perez - “A mathematical formalization of the parallel replica dynamics”, Monte Carlo Methods Appl. 18 (2012) no. 2, p. 119-146 | Zbl | DOI | MR
\bibitem{le-peutrec-10}D. Le Peutrec - “Small eigenvalues of the Neumann realization of the semiclassical Witten Laplacian”, Ann. Fac. Sci. Toulouse Math. (6) 19 (2010) no. 3-4, p. 735-809 | DOI | MR | Numdam | Zbl
\bibitem{dorianmichel}D. Le Peutrec \& L. Michel - “Sharp spectral asymptotics for non-reversible metastable diffusion processes”, Probab. Math. Phys. 1 (2020) no. 1, p. 3-53 | DOI | MR | Zbl
\bibitem{DoNe2}D. Le Peutrec \& B. Nectoux - “Small eigenvalues of the Witten Laplacian with Dirichlet boundary conditions: the case with critical points on the boundary”, Anal. PDE 14 (2021) no. 8, p. 2595-2651 | DOI | MR | Zbl
\bibitem{le-peutrec_nier_viterbo-13}D. Le Peutrec, F. Nier \& C. Viterbo - “Precise Arrhenius law for $p$-forms: the Witten Laplacian and Morse-Barannikov complex”, Ann. Henri Poincaré 14 (2013) no. 3, p. 567-610 | DOI | MR | Zbl
\bibitem{peutrec2020bar}D. Le Peutrec, F. Nier \& C. Viterbo - Bar codes of persistent cohomology and Arrhenius law for $p$-forms, Astérisque, vol. 450, Société Mathématique de France, Paris, 2024 | DOI | Zbl
\bibitem{lee2020non}J. Lee \& I. Seo - “Non-reversible metastable diffusions with Gibbs invariant measure I: Eyring–Kramers formula”, Probab. Theory Related Fields 182 (2022) no. 3, p. 849-903 | DOI | MR | Zbl
\bibitem{lelievre-18}T. Lelièvre - “Mathematical foundations of accelerated molecular dynamics methods”, in Handbook of materials modeling (W. Andreoni \& S. Yip, eds.), 2018 | DOI
\bibitem{IHPLLN}T. Lelièvre, D. Le Peutrec \& B. Nectoux - “Exit event from a metastable state and Eyring-Kramers law for the overdamped Langevin dynamics”, in Stochastic dynamics out of equilibrium, Springer Proc. Math. Stat., vol. 282, Springer, Cham, 2019, p. 331-363 | DOI | Zbl
\bibitem{DLLN-saddle2}T. Lelièvre, D. Le Peutrec \& B. Nectoux - “The exit from a metastable state: concentration of the exit point distribution on the low energy saddle points, part 2”, Stochastic Partial Differ. Equ. Anal. Comput. 10 (2022), p. 317-357 | MR | Zbl
\bibitem{LeNi}T. Lelièvre \& F. Nier - “Low temperature asymptotics for quasistationary distributions in a bounded domain”, Anal. PDE 8 (2015) no. 3, p. 561-628 | DOI | MR | Zbl
\bibitem{mathieu1995spectra}P. Mathieu - “Spectra, exit times and long time asymptotics in the zero-white-noise limit”, Stochastics 55 (1995) no. 1-2, p. 1-20 | Zbl
\bibitem{MS77}B. J. Matkowsky \& Z. Schuss - “The exit problem for randomly perturbed dynamical systems”, SIAM J. Appl. Math. 33 (1977) no. 2, p. 365-382 | DOI | MR | Zbl
\bibitem{MaSc}B. J. Matkowsky \& Z. Schuss - “The exit problem: a new approach to diffusion across potential barriers”, SIAM J. Appl. Math. 36 (1979) no. 3, p. 604-623 | MR | Zbl
\bibitem{menz-schlichting-14}G. Menz \& A. Schlichting - “Poincaré and logarithmic Sobolev inequalities by decomposition of the energy landscape”, Ann. Probab. 42 (2014) no. 5, p. 1809-1884 | Zbl | MR | DOI
\bibitem{michel2017small}L. Michel - “About small eigenvalues of Witten Laplacian”, Pure Appl. Anal. 1 (2019) no. 2, p. 149-206 | Zbl | DOI | MR
\bibitem{miclo-95}L. Miclo - “Comportement de spectres d’opérateurs de Schrödinger à basse température”, Bull. Sci. Math. 119 (1995) no. 6, p. 529-554 | MR | Zbl
\bibitem{milnor1997topology}J. W. Milnor - Topology from the differentiable viewpoint, Princeton Landmarks in Math., Princeton University Press, Princeton, NJ, 1997 | MR | Zbl
\bibitem{PhysRevLett87126101}F. Montalenti, M. R. Sørensen \& A. F. Voter - “Closing the gap between experiment and theory: crystal growth by temperature accelerated dynamics”, Phys. Rev. Lett. 87 (2001), article ID 126101, 4 pages | DOI
\bibitem{schuss90}T. Naeh, M. M. Klosek, B. J. Matkowsky \& Z. Schuss - “A direct approach to the exit problem”, SIAM J. Appl. Math. 50 (1990) no. 2, p. 595-627 | DOI | MR | Zbl
\bibitem{NePC}B. Nectoux - “Mean exit time for the overdamped Langevin process: the case with critical points on the boundary”, Comm. Partial Differential Equations 46 (2021) no. 9, p. 1789-1829 | DOI | Zbl
\bibitem{pande2010everything}V. S. Pande, K. Beauchamp \& G. R. Bowman - “Everything you wanted to know about Markov state models but were afraid to ask”, Methods 52 (2010) no. 1, p. 99-105 | DOI
\bibitem{lelievre-perez-24}D. Perez \& T. Lelièvre - “Recent advances in accelerated molecular dynamics methods: Theory and applications”, in Comprehensive Computational Chemistry (M. Yáñez \& R. J. Boyd, eds.), vol. 3, Elsevier, 2024, p. 360-383 | DOI
\bibitem{perthame-90}B. Perthame - “Perturbed dynamical systems with an attracting singularity and weak viscosity limits in Hamilton-Jacobi equations”, Trans. Amer. Math. Soc. 317 (1990) no. 2, p. 723-748 | DOI | MR | Zbl
\bibitem{schuette-98}C. Schütte - Conformational dynamics: modelling, theory, algorithm and application to biomolecules, Habilitation dissertation, Free University Berlin, 1998
\bibitem{schuette-sarich-13}C. Schütte \& M. Sarich - Metastability and Markov state models in molecular dynamics, Courant Lecture Notes in Math., vol. 24, American Mathematical Society, Providence, RI, 2013 | DOI | Zbl | MR
\bibitem{GSchw}G. Schwarz - Hodge decomposition—a method for solving boundary value problems, Lect. Notes in Math., vol. 1607, Springer-Verlag, Berlin, 1995 | Zbl | DOI | MR
\bibitem{simon1983semiclassical}B. Simon - “Semiclassical analysis of low lying eigenvalues. I. Non-degenerate minima: Asymptotic expansions”, Ann. Inst. H. Poincaré Phys. Théor. 38 (1983) no. 3, p. 295-308 | Numdam | Zbl
\bibitem{sorensen-voter-00}M. R. Sorensen \& A. F. Voter - “Temperature-accelerated dynamics for simulation of infrequent events”, J. Chem. Phys. 112 (2000) no. 21, p. 9599-9606 | DOI
\bibitem{PhysRevB66205415}J. A. Sprague, F. Montalenti, B. P. Uberuaga, J. D. Kress \& A. F. Voter - “Simulation of growth of Cu on Ag(001) at experimental deposition rates”, Phys. Rev. B 66 (2002), article ID 205415, 10 pages | DOI
\bibitem{sugiura1996}M. Sugiura - “Exponential asymptotics in the small parameter exit problem”, Nagoya Math. J. 144 (1996), p. 137-154 | DOI | MR | Zbl
\bibitem{sugiura2001}M. Sugiura - “Asymptotic behaviors on the small parameter exit problems and the singularly perturbation problems”, Ryukyu Math. J. 14 (2001), p. 79-118 | MR | Zbl
\bibitem{thom1954quelques}R. Thom - “Quelques propriétés globales des variétés différentiables”, Comment. Math. Helv. 28 (1954) no. 1, p. 17-86 | DOI | Zbl | MR
\bibitem{tiwary2013metadynamics}P. Tiwary \& M. Parrinello - “From metadynamics to dynamics”, Phys. Rev. Lett. 111 (2013), article ID 230602, 5 pages | DOI
\bibitem{tong1952some}H. Tong - “Some characterizations of normal and perfectly normal spaces”, Duke Math. J. 19 (1952) no. 2, p. 289-292 | DOI | MR | Zbl
\bibitem{voter-97}A. F. Voter - “A method for accelerating the molecular dynamics simulation of infrequent events”, J. Chem. Phys. 106 (1997) no. 11, p. 4665-4677 | DOI
\bibitem{voter-05}A. F. Voter - “Introduction to the kinetic Monte Carlo method”, Radiation Effects in Solids, Springer, NATO Publishing Unit (2005)
\bibitem{wales-03}D. J. Wales - Energy landscapes: applications to clusters, biomolecules and glasses, Cambridge Molecular Science, Cambridge University Press, Cambridge, 2004 | DOI
\bibitem{witten-82}E. Witten - “Supersymmetry and Morse theory”, J. Differential Geom. 17 (1982) no. 4, p. 661-692 | Zbl | MR
\end{thebibliography}
